\documentclass[10pt,a4paper]{amsart}
\usepackage[margin=19mm]{geometry}
\usepackage{amsmath,amssymb,mathtools}
\usepackage[scr=rsfs,cal=euler]{mathalfa}
\usepackage{xfrac}
\usepackage[unicode,hidelinks,bookmarks=false]{hyperref}
\newcommand{\ie}{i.e.\ }
\newcommand{\cf}{cf.\ }
\newcommand{\resp}{respectively\ }
\newcommand{\Subsection}{\subsection}
\providecommand{\nobreakdash}{\nobreak\hbox{-}}
\setlength{\emergencystretch}{2em}
\multlinegap0pt
\usepackage{bm}
\usepackage{bbm}
\usepackage{dsfont}
\usepackage{xspace}
\usepackage{cancel}
\usepackage{mathtools}
\usepackage{amsthm}
\makeatletter
\@ifundefined{thm}{\newtheorem{thm}{Thm}}{}
\@ifundefined{lem}{\newtheorem{lem}[thm]{Lemma}}{}
\makeatother
\makeatletter
\expandafter\ifx\csname abstract\endcsname\relax
\renewenvironment{abstract}{
  \begin{center}
  \vspace{-0.5em}
  \rule{\textwidth}{0.4pt}\\[4pt]
  \begin{minipage}{0.95\textwidth}
  \small
}{
  \end{minipage}\\[4pt]
  \rule{\textwidth}{0.4pt}
  \end{center}
  \vspace{1em}
}
\fi
\makeatother
\title[Kinetic Theory and the Mechanics of Isothermal Gas Spheres: ]{Kinetic Theory and the Mechanics of Isothermal Gas Spheres: Derivation of the Classical Emden--Chandrasekhar Equation via the Vlasov--Poisson Formalism}
\author{Steven D Miller}
\date{Source: arXiv:2511.08814v1 (CC BY 4.0)}
\begin{document}
\maketitle

\noindent\emph{Source and licence.} Kinetic Theory and the Mechanics of Isothermal Gas Spheres: Derivation of the Classical Emden--Chandrasekhar Equation via the Vlasov--Poisson Formalism. Version arXiv:2511.08814v1 (CC BY 4.0), distributed under the Creative Commons Attribution 4.0 International licence, as recorded in the arXiv licence metadata for this exact version and shown on the abstract page https://arxiv.org/abs/2511.08814v1. Full rights notice: licenses/arxiv-2511.08814-CC-BY-4.0.txt.\par\smallskip

\noindent\textit{Editorial scope.} Counted unit: the Lem whose source label is \texttt{None} in the article (source0.tex lines 1410--1431, source label \texttt{None}), delivered together with the article's own complete original proof of it (source0.tex lines 1432--1534, one \texttt{proof} environment of 103 lines). No mathematics is written, repaired or re-derived: statement and proof are the article's own text line for line. Required context retained uncounted: none. Every definition, display and label of those lines is retained unchanged. Omitted from this package and not redistributed: 1--1409, 1535--1911. Nothing inside a retained excerpt is omitted or altered: every retained body is delivered exactly as the article's lines are written, so no omission receipt of any kind accompanies this package. Citations: the retained excerpts carry no citation command at all, so no bibliography entry of the article is transcribed and no reference is redistributed. Reference adaptation: none was needed and none was made; every reference inside the delivered unit resolves inside it, so no unresolved reference or citation is produced. Printed slips kept literal: no printed word, symbol, hypothesis or number of the counted statement or of its proof was altered. Layout and class: the article's own class is \texttt{article} with packages including inputenc, amsmath, amssymb, amsthm, amsfonts, amsbsy, mathtools, mathrsfs, bm, bbm, stmaryrd, relsize, graphicx, tikz; the wrapper is the lane's pinned \texttt{amsart} editorial wrapper at 10pt on A4, which restates the theorem-like environments the retained text needs and adds the article's own macro definitions that the retained text needs (none), with extra wrapper packages (bm, bbm, dsfont, xspace, cancel, mathtools). The delivered numbering is the unit's own. Assets: no retained span contains a figure environment, a graphics-inclusion command, a PDF-page inclusion, a table or a TikZ picture; no image file is copied into this package, the package declares no asset dependency and no image file is redistributed with it. Licence: version arXiv:2511.08814v1 of this article is distributed under the Creative Commons Attribution 4.0 International licence, as recorded in the arXiv licence metadata for this exact version and shown on the abstract page https://arxiv.org/abs/2511.08814v1. A full rights notice is delivered at licenses/arxiv-2511.08814-CC-BY-4.0.txt. Characters: every retained excerpt is pure ASCII, so the delivered engine, XeLaTeX with its default Latin Modern text font, reports no missing character.
\begin{lem}[\textbf{Scalar virial theorem}]
Let $\mathcal{F}(\mathbf{x},\mathbf{p})$ be a stationary solution of the Vlasov--Poisson equation
\begin{equation}
\frac{\mathbf{p}}{m}\!\cdot\!\nabla_{\mathbf{x}} f - m\nabla_{\mathbf{x}}\Phi(\mathbf{x})\!\cdot\!\nabla_{\mathbf{p}} f = 0
\qquad\text{on }\Gamma=\Omega\times\mathbb{R}^3,
\end{equation}
where $\Omega\subset\mathbb{R}^3$ is the physical domain (either all space or a bounded domain). Assume the decay / boundary conditions
\begin{itemize}
  \item $\mathcal{F}(\mathbf{x},\mathbf{p})\to 0$ sufficiently fast as $|\mathbf{p}|\to\infty$ so that momentum-space surface integrals vanish;
  \item either $\Omega$ is bounded and $f|_{\partial\Omega}=0$, or $\Omega=\mathbb{R}^3$ and $f,\Phi\to 0$ sufficiently rapidly as $|\mathbf{x}|\to\infty$ so that spatial surface integrals vanish.
\end{itemize}
Then, with
\begin{equation}
\langle {K}\rangle = \int_{\Gamma}\frac{|\mathbf{p}|^2}{2m}\,\mathcal{F}(\mathbf{x},\mathbf{p})\,d^3x\,d^3p,
\qquad
U = \tfrac12\int_{\Omega}\rho(\mathbf{x})\Phi(\mathbf{x})\,d^3x,
\end{equation}
one has the scalar virial identity
\begin{equation}
2\langle {K}\rangle + U = 0.
\end{equation}
\end{lem}

\begin{proof}
Start from the stationary Vlasov equation and multiply by the scalar $(\mathbf{x}\!\cdot\!\mathbf{p})$, then integrate over phase space:
\begin{equation}
0
= \int_{\Gamma} (\mathbf{x}\!\cdot\!\mathbf{p})\Big[\frac{\mathbf{p}}{m}\!\cdot\!\nabla_{\mathbf{x}} f - m\nabla_{\mathbf{x}}\Phi\!\cdot\!\nabla_{\mathbf{p}} f\Big]\,d^3x\,d^3p
= \mathscr{B} + \mathscr{C},
\end{equation}
where we define
\begin{equation}
\mathscr{B} := \int_{\Gamma} (\mathbf{x}\!\cdot\!\mathbf{p})\frac{\mathbf{p}}{m}\!\cdot\!\nabla_{\mathbf{x}} f\,d^3x\,d^3p,
\qquad
\mathscr{C} := -\int_{\Gamma} (\mathbf{x}\!\cdot\!\mathbf{p})\,m\nabla_{\mathbf{x}}\Phi\!\cdot\!\nabla_{\mathbf{p}} f\,d^3x\,d^3p.
\end{equation}

\paragraph{Integration by parts in \(\mathbf{x}\) for \(\mathscr{B}\).}
Write componentwise (Einstein summation):
\begin{equation}
\mathscr{B} = \frac{1}{m}\int_{\Gamma} x_i p_i p_j \frac{\partial f}{\partial x_j}\,d^3x\,d^3p.
\end{equation}
Observe the identity
\begin{equation}
x_i p_i p_j \frac{\partial f}{\partial x_j}
= \frac{\partial}{\partial x_j}\big(x_i p_i p_j f\big) - p_j p_j f
= \nabla_{\mathbf{x}}\!\cdot\!\big( (\mathbf{x}\!\cdot\!\mathbf{p})\,\mathbf{p}\, f \big) - |\mathbf{p}|^2 f.
\end{equation}
Integrate over \(\Omega\) and then over \(\mathbf{p}\). The divergence term produces a surface integral on \(\partial\Omega\) which vanishes under the hypotheses, hence
\begin{equation}
\mathscr{B} = -\frac{1}{m}\int_{\Gamma} |\mathbf{p}|^2 f\,d^3x\,d^3p
= -\int_{\Gamma}\frac{|\mathbf{p}|^2}{m} f\,d^3x\,d^3p
= -2\langle {K}\rangle.
\end{equation}

\paragraph{Integration by parts in \(\mathbf{p}\) for \(\mathscr{C}\).}
For fixed \(\mathbf{x}\), set \(\mathbf{a}(\mathbf{x}):=\nabla_{\mathbf{x}}\Phi(\mathbf{x})\) (independent of \(\mathbf{p}\)). Consider the inner momentum integral
\begin{equation}
\mathscr{J}(\mathbf{x}) := \int_{\mathbb{R}^3} (\mathbf{x}\!\cdot\!\mathbf{p})\,\big(\mathbf{a}(\mathbf{x})\!\cdot\!\nabla_{\mathbf{p}} \mathcal{F}(\mathbf{x},\mathbf{p})\big)\,d^3p.
\end{equation}
Integrate by parts in \(\mathbf{p}\); the boundary term at $|\mathbf{p}|\to\infty$ vanishes by decay of \(f\). Using \(\nabla_{\mathbf{p}}(\mathbf{x}\!\cdot\!\mathbf{p})=\mathbf{x}\) we get
\begin{equation}
\mathscr{J}(\mathbf{x}) = -\int_{\mathbb{R}^3} \mathcal{F}(\mathbf{x},\mathbf{p})\,\mathbf{a}(\mathbf{x})\!\cdot\!\mathbf{x}\,d^3p
= -(\mathbf{x}\!\cdot\!\nabla_{\mathbf{x}}\Phi(\mathbf{x}))\int_{\mathbb{R}^3} \mathcal{F}(\mathbf{x},\mathbf{p})\,d^3p
= -(\mathbf{x}\!\cdot\!\nabla_{\mathbf{x}}\Phi(\mathbf{x}))\,\rho(\mathbf{x}).
\end{equation}
Therefore, recalling the outer factor \(-m\) in the definition of \(\mathscr{C}\),
\begin{equation}
\mathscr{C} = -m\int_{\Omega} \mathscr{J}(\mathbf{x})\,d^3x
= -m\int_{\Omega} \big[-(\mathbf{x}\!\cdot\!\nabla\Phi)\rho(\mathbf{x})\big]\,d^3x
= +m\int_{\Omega} \rho(\mathbf{x})\,\mathbf{x}\!\cdot\!\nabla\Phi(\mathbf{x})\,d^3x.
\end{equation}
With our convention that $f$ is the \emph{mass} density in phase space (so that $\rho=\int f\,d^3p$ is already the mass density), the extra $m$ is redundant and cancels; hence one obtains simply
\begin{equation}
\mathscr{C} \;=\; \int_{\Omega} \rho(\mathbf{x})\,\mathbf{x}\!\cdot\!\nabla\Phi(\mathbf{x})\,d^3x \;=: \; \mathcal{W}.
\end{equation}

\paragraph{Combine the two pieces.}
From \(\mathscr{B}+\mathscr{C}=0\) we have
\begin{equation}
-2\langle {K}\rangle + \mathcal{W} = 0 \qquad\Longrightarrow\qquad 2\langle {K}\rangle = \mathcal{W}.
\label{eq:virial_intermediate}
\end{equation}

\paragraph{Auxiliary identity: \(\mathcal{W} = -\dfrac{1}{2}\!\displaystyle\int_{\Omega}\rho\Phi\,d^3x\).}
Using Poisson's equation \(\Delta\Phi=4\pi G\rho\),
\begin{equation}
\mathcal{W} = \frac{1}{4\pi G}\int_{\Omega} (\Delta\Phi)\,(\mathbf{x}\!\cdot\!\nabla\Phi)\,d^3x.
\end{equation}
Set $a(\mathbf{x}):=\mathbf{x}\!\cdot\!\nabla\Phi(\mathbf{x})$. Integrating by parts and dropping spatial surface terms gives
\begin{equation}
\int_{\Omega} (\Delta\Phi)\,a\,d^3x = -\int_{\Omega} \nabla a\!\cdot\!\nabla\Phi\,d^3x.
\end{equation}
A short calculation yields
\begin{equation}
\nabla a\!\cdot\!\nabla\Phi = |\nabla\Phi|^2 + \tfrac12\,\mathbf{x}\!\cdot\!\nabla\big(|\nabla\Phi|^2\big).
\end{equation}
Integrating the second term by parts (and discarding the boundary term) gives
\begin{equation}
\int_{\Omega} \mathbf{x}\!\cdot\!\nabla\big(|\nabla\Phi|^2\big)\,d^3x = -3\int_{\Omega} |\nabla\Phi|^2\,d^3x,
\end{equation}
hence
\begin{equation}
\int_{\Omega} \nabla a\!\cdot\!\nabla\Phi\,d^3x = \tfrac12\int_{\Omega} |\nabla\Phi|^2\,d^3x.
\end{equation}
Consequently
\begin{equation}
\mathcal{W} = -\frac{1}{4\pi G}\cdot \tfrac12\int_{\Omega} |\nabla\Phi|^2\,d^3x
= -\frac{1}{8\pi G}\int_{\Omega} |\nabla\Phi|^2\,d^3x.
\end{equation}
Multiplying Poisson's equation by \(\Phi\) and integrating by parts gives
\begin{equation}
\int_{\Omega} \rho\,\Phi\,d^3x = -\frac{1}{4\pi G}\int_{\Omega} |\nabla\Phi|^2\,d^3x.
\end{equation}
Comparing the two relations yields
\begin{equation}
\mathcal{W} = -\tfrac12 \int_{\Omega} \rho\,\Phi\,d^3x = -\,U.
\end{equation}

\paragraph{Final virial identity.}
Insert this into \eqref{eq:virial_intermediate}:
\begin{equation}
2\langle {K}\rangle = \mathcal{W} = -U \quad\Longrightarrow\quad 2\langle {K}\rangle + U = 0,
\end{equation}
which is the usual scalar virial theorem in astrophysics.
\end{proof}

\end{document}
